\chapter{De stelling van Thales}\label{ch:g9:thales}

Hoe meet je de hoogte van een \omterm{def:g8:solids:def}{piramide} zonder erop te klimmen? Thales van Milete
vergeleek haar schaduw met de schaduw van een stok. Zijn stelling --- \omterm{def:g6:lines:perp}{parallelle
lijnen} verdelen lijnstukken in evenredige stukken --- is het wiskundige hart van
maquettes, kaarten en vergrotingen, en een van de oudste stellingen met een naam.

\section{De stelling}

\begin{theorem}[Thales]\label{thm:g9:thales:direct}
Zij twee \omterm{def:g6:lines:objects}{lijnen} die elkaar in een punt $A$ snijden. Neem $M$ op de eerste \omterm{def:g6:lines:objects}{lijn} en
$N$ op de tweede, met $B$ verder op de eerste en $C$ verder op de tweede. Zijn de
\omterm{def:g6:lines:objects}{lijnen} $(MN)$ en $(BC)$ \emph{\omterm{def:g6:lines:perp}{parallel}}, dan hebben de driehoeken $AMN$ en $ABC$
evenredige zijden:
\[
\frac{AM}{AB} = \frac{AN}{AC} = \frac{MN}{BC}.
\]
\end{theorem}
\admitted

\begin{omfigure}
\begin{tikzpicture}[scale=1.05]
  \coordinate (A) at (0,0);
  \coordinate (B) at (5.0,0.6);
  \coordinate (C) at (3.2,2.8);
  \coordinate (M) at ($(A)!0.55!(B)$);
  \coordinate (N) at ($(A)!0.55!(C)$);
  \draw[thick] (A) -- (B) -- (C) -- cycle;
  \draw[omThm, very thick] (M) -- (N);
  \draw[omThm, very thick] (B) -- (C);
  \fill (A) circle (1.5pt) node[left, font=\small] {$A$};
  \fill (B) circle (1.5pt) node[right, font=\small] {$B$};
  \fill (C) circle (1.5pt) node[above, font=\small] {$C$};
  \fill (M) circle (1.5pt) node[below, font=\small] {$M$};
  \fill (N) circle (1.5pt) node[above left, font=\small] {$N$};
\end{tikzpicture}\hspace{1.1cm}%
\begin{tikzpicture}[scale=1.05]
  \coordinate (A) at (0,0);
  \coordinate (B) at (2.6,-1.3);
  \coordinate (C) at (1.4,-2.2);
  \coordinate (M) at ($(A)!-0.62!(B)$);
  \coordinate (N) at ($(A)!-0.62!(C)$);
  \draw[thick] (M) -- (B);
  \draw[thick] (N) -- (C);
  \draw[omThm, very thick] (M) -- (N);
  \draw[omThm, very thick] (B) -- (C);
  \fill (A) circle (1.5pt) node[right=2pt, font=\small] {$A$};
  \fill (B) circle (1.5pt) node[right, font=\small] {$B$};
  \fill (C) circle (1.5pt) node[below, font=\small] {$C$};
  \fill (M) circle (1.5pt) node[left, font=\small] {$M$};
  \fill (N) circle (1.5pt) node[above, font=\small] {$N$};
\end{tikzpicture}

{\small De twee standen van Thales: geneste driehoeken (links) en de
``vlinderstand'', waarbij $M$ en $N$ aan de andere kant van $A$ liggen (rechts).
In beide is $(MN) \parallel (BC)$ en zijn de drie verhoudingen gelijk.}
\end{omfigure}

\begin{method}[Een lengte berekenen met Thales]\label{met:g9:thales:compute}
\begin{enumerate}
  \item Bepaal het punt $A$ waar de twee \omterm{def:g6:lines:objects}{lijnen} elkaar snijden en de twee
        \omterm{def:g6:lines:perp}{parallelle lijnen}; benoem de twee driehoeken;
  \item schrijf de drie gelijke verhoudingen op, telkens ``kleine driehoek gedeeld
        door grote driehoek'';
  \item houd de gelijkheid van de twee verhoudingen over waarin de drie bekende
        lengten en de onbekende voorkomen;
  \item los de verhouding die je zo krijgt op door kruislings te vermenigvuldigen.
\end{enumerate}
\end{method}

\begin{example}\label{ex:g9:thales:compute}
Stel in de linkerstand hierboven dat $AM = 3$, $AB = 5$, $AN = 2.4$ en
$MN = 3.3$, en laten we $AC$ en $BC$ berekenen. Thales geeft
\[
\frac{AM}{AB} = \frac{AN}{AC} = \frac{MN}{BC},
\qquad\text{dat wil zeggen}\qquad
\frac35 = \frac{2.4}{AC} = \frac{3.3}{BC}.
\]
Uit $\frac35 = \frac{2.4}{AC}$ volgt door kruislings te vermenigvuldigen:
$3 \times AC = 5 \times 2.4 = 12$, dus $AC = 4$. Uit $\frac35 = \frac{3.3}{BC}$:
$3 \times BC = 16.5$, dus $BC = 5.5$.
\end{example}

\begin{example}[Thales in het echte leven]\label{ex:g9:thales:stick}
Een verticale stok van $1$ m hoog werpt een schaduw van $1.5$ m, terwijl een boom
op hetzelfde ogenblik een schaduw van $12$ m werpt. Zonnestralen zijn \omterm{def:g6:lines:perp}{parallel},
dus staan de driehoeken van stok-met-schaduw en boom-met-schaduw in de stand van
Thales:
\[
\frac{\text{hoogte van de boom}}{1} = \frac{12}{1.5},
\qquad\text{dus is de boom } 8 \text{ m hoog.}
\]
\end{example}

\section{De omgekeerde stelling}

\begin{theorem}[Omgekeerde stelling van Thales]\label{thm:g9:thales:converse}
Met de punten geplaatst als in \cref{thm:g9:thales:direct} ($M$ op $(AB)$, $N$ op
$(AC)$, in dezelfde orde op beide \omterm{def:g6:lines:objects}{lijnen}): geldt
\[
\frac{AM}{AB} = \frac{AN}{AC},
\]
dan zijn de \omterm{def:g6:lines:objects}{lijnen} $(MN)$ en $(BC)$ \omterm{def:g6:lines:perp}{parallel}.
\end{theorem}
\admitted

\begin{method}[Parallellisme bewijzen of weerleggen]\label{met:g9:thales:converse}
\begin{enumerate}
  \item Bereken de twee verhoudingen $\dfrac{AM}{AB}$ en $\dfrac{AN}{AC}$
        afzonderlijk (als \omterm{def:g9:fractions:fraction}{breuken}, niet als \omterm{def:g6:decimals:rounding}{afrondingen});
  \item zijn ze gelijk \emph{en} liggen de punten in dezelfde orde op de twee
        \omterm{def:g6:lines:objects}{lijnen}, dan zijn de \omterm{def:g6:lines:objects}{lijnen} \omterm{def:g6:lines:perp}{parallel} (omgekeerde stelling van Thales);
  \item verschillen ze, dan zijn de \omterm{def:g6:lines:objects}{lijnen} \emph{niet} \omterm{def:g6:lines:perp}{parallel} --- want waren ze
        dat, dan zou de stelling van Thales gelijke verhoudingen opleggen.
\end{enumerate}
\end{method}

\begin{example}\label{ex:g9:thales:converse}
Op \omterm{def:g6:lines:objects}{lijn} 1 is $AM = 4$ en $AB = 10$; op \omterm{def:g6:lines:objects}{lijn} 2 is $AN = 6$ en $AC = 15$. Dan is
\[
\frac{AM}{AB} = \frac{4}{10} = \frac25,
\qquad
\frac{AN}{AC} = \frac{6}{15} = \frac25 :
\]
gelijke verhoudingen, dezelfde orde: $(MN) \parallel (BC)$.

Met $AC = 14$ in de plaats: $\frac{6}{14} = \frac37 \neq \frac25$, dus zijn de
\omterm{def:g6:lines:objects}{lijnen} niet \omterm{def:g6:lines:perp}{parallel}.
\end{example}

\section{Vergroten en verkleinen}

\begin{definition}[Een figuur schalen]\label{def:g9:thales:scaling}
Een figuur vergroten of verkleinen met de \emph{schaalfactor}\index{schaalfactor}
$k > 0$ betekent \emph{al} haar lengten met $k$ vermenigvuldigen: een
\emph{vergroting} wanneer $k > 1$, een \emph{verkleining} wanneer $k < 1$. In de
stand van Thales is de driehoek $AMN$ de verkleining van $ABC$ met de factor
$k = \frac{AM}{AB}$.
\end{definition}

\begin{proposition}[Gevolg voor hoeken en oppervlakten]\label{prop:g9:thales:scaling}
Schalen met de factor $k$ behoudt de hoeken en de vorm van figuren,
vermenigvuldigt elke lengte met $k$, en vermenigvuldigt elke \emph{\omterm{def:g6:measure:area}{oppervlakte}}
met $k^2$.
\end{proposition}
\admitted

\begin{omfigure}
\begin{tikzpicture}[scale=0.72]
  \draw[thick, omDef, fill=omDef!8] (0,0) rectangle (2,1.2);
  \draw[thick, omThm, fill=omThm!8] (3.4,0) rectangle (7.4,2.4);
  \node[font=\small, omDef] at (1,-0.45) {$2 \times 1.2$};
  \node[font=\small, omThm] at (5.4,-0.45) {$4 \times 2.4$};
  \node[font=\small] at (1,0.6) {opp.\ $2.4$};
  \node[font=\small] at (5.4,1.2) {opp.\ $9.6$};
\end{tikzpicture}

{\small De lengten verdubbelen ($k = 2$) vermenigvuldigt de \omterm{def:g6:measure:area}{oppervlakte} met
$k^2 = 4$: vier exemplaren van de kleine rechthoek betegelen de grote.}
\end{omfigure}

\begin{example}
Een foto van $10 \times 15$ cm wordt vergroot met \omterm{def:g9:thales:scaling}{schaalfactor} $k = 3$: de afdruk
meet $30 \times 45$ cm. Haar \omterm{def:g6:measure:area}{oppervlakte} gaat van $150$ cm$^2$ naar
$150 \times 9 = 1350$ cm$^2$ --- $9$ keer groter, niet $3$.
\end{example}

\section{Oefeningen}

\begin{exercise}[$\star$]\label{exo:g9:thales:1}
De \omterm{def:g6:lines:objects}{lijnen} $(MN)$ en $(BC)$ zijn \omterm{def:g6:lines:perp}{parallel}, met $M$ op $[AB]$ en $N$ op $[AC]$; en
$AM = 2$, $AB = 6$, $AN = 3$, $BC = 9$. Bereken $AC$ en $MN$.
\end{exercise}

\begin{exercise}[$\star$]\label{exo:g9:thales:2}
Dezelfde stand: $AM = 5$, $MB = 3$ (let op: $MB$, niet $AB$!), $AN = 4$. Bereken
$AC$.
\end{exercise}

\begin{exercise}[$\star$]\label{exo:g9:thales:3}
In een vlinderstand ligt $A$ tussen $M$ en $B$ en tussen $N$ en $C$, met
$(MN) \parallel (BC)$; verder is $AM = 3$, $AB = 7.5$, $AN = 2$, $MN = 2.6$.
Bereken $AC$ en $BC$.
\end{exercise}

\begin{exercise}[$\star$]\label{exo:g9:thales:4}
$M$ ligt op $[AB]$ met $AM = 6$, $AB = 8$; $N$ ligt op $[AC]$ met $AN = 9$,
$AC = 12$. Zijn de \omterm{def:g6:lines:objects}{lijnen} $(MN)$ en $(BC)$ \omterm{def:g6:lines:perp}{parallel}?
\end{exercise}

\begin{exercise}[$\star\star$]\label{exo:g9:thales:5}
Zoals hierboven, met $AM = 4$, $AB = 6$, $AN = 5$, $AC = 8$. Zijn $(MN)$ en
$(BC)$ \omterm{def:g6:lines:perp}{parallel}? Verantwoord zorgvuldig.
\end{exercise}

\begin{exercise}[$\star\star$]\label{exo:g9:thales:6}
Een persoon van $1.8$ m staat op $2$ m van de voet van een straatlantaarn; haar
schaduw meet $3$ m. Teken de stand van Thales die de lantaarn, de persoon en het
uiteinde van de schaduw vormen, en bereken de hoogte van de lantaarn.
\end{exercise}

\begin{exercise}[$\star\star$]\label{exo:g9:thales:7}
Een kaart heeft schaal $1 : 25\,000$ (de lengten op de kaart zijn de echte lengten
vermenigvuldigd met $k = \frac{1}{25000}$).
\begin{enumerate}
  \item Twee dorpen liggen op de kaart $6.8$ cm van elkaar. Wat is de echte
        afstand, in km?
  \item Een bos heeft op de kaart een \omterm{def:g6:measure:area}{oppervlakte} van $8$ cm$^2$. Wat is zijn echte
        \omterm{def:g6:measure:area}{oppervlakte}, in km$^2$?
\end{enumerate}
\end{exercise}

\begin{exercise}[$\star\star$]\label{exo:g9:thales:8}
De driehoek $ABC$ heeft $AB = 12$, $AC = 15$, $BC = 18$. Het punt $M$ op $[AB]$
voldoet aan $AM = 8$, en de \omterm{def:g6:lines:objects}{lijn} door $M$ \omterm{def:g6:lines:perp}{parallel} met $(BC)$ snijdt $[AC]$ in
$N$.
\begin{enumerate}
  \item Bereken $AN$ en $MN$.
  \item Wat is de \omterm{def:g9:thales:scaling}{schaalfactor} van $ABC$ naar $AMN$? Vergelijk de omtrekken van de
        twee driehoeken.
\end{enumerate}
\end{exercise}

\begin{exercise}[$\star\star\star$]\label{exo:g9:thales:9}
Zij $ABCD$ een trapezium met $(AB) \parallel (CD)$, $AB = 4$ en $CD = 6$, waarvan
de diagonalen $[AC]$ en $[BD]$ elkaar in $O$ snijden. Bereken met Thales in de
vlinderstand rond $O$ de verhouding $\dfrac{OA}{OC}$, en toon aan dat $O$ beide
diagonalen in dezelfde verhouding verdeelt.
\end{exercise}

\section{Opgave: de lat zonder maatstreep, de camera obscura en de drie
gemiddelden van een trapezium}

\begin{problem}[{Weekendopgave --- Thales aan het werk: elk lijnstuk in gelijke
delen verdelen, de Zon meten met een schoenendoos, en een trapezium waarin alle
drie de beroemde gemiddelden samenkomen}]
\label{pb:g9:thales:1}
De stelling van Thales is de wiskunde van de \emph{stralen}: zonnestralen,
lichtstralen door een gaatje, potloodstralen uit een \omterm{def:g5:solids:def}{hoekpunt}. Deze opgave gebruikt
haar op drie manieren --- als constructiegereedschap (een \omterm{def:g6:lines:objects}{lijnstuk} van
\emph{willekeurige} lengte, zelfs $\sqrt2$, in volmaakt gelijke delen verdelen),
als meetinstrument (de \omterm{def:g4:geometry:circle}{diameter} van de Zon, met een schoenendoos en een munt), en
als vergrootglas op het trapezium van \cref{exo:g9:thales:9}, waar het
rekenkundig, het meetkundig en het \omterm{pb:g8:speed:1}{harmonisch gemiddelde} van deze reeks alle in
één figuur opduiken.

\medskip
\noindent\textbf{Deel I --- Verdelen met een lat zonder maatstreep.}
\begin{enumerate}
  \item Teken een \omterm{def:g6:lines:objects}{lijnstuk} $[AB]$ van $7$ cm. Om het met passer en een lat zonder
        maatstreep in drie gelijke delen te verdelen: teken een willekeurige straal
        uit $A$ (niet door $B$); zet er drie gelijke passerlengten op uit, wat de
        punten $P_1$, $P_2$, $P_3$ geeft; verbind $P_3$ met $B$; teken de
        parallellen met $(P_3 B)$ door $P_1$ en $P_2$. Voer de constructie uit en
        markeer waar de parallellen $[AB]$ snijden.
  \item Verantwoord met \cref{thm:g9:thales:direct} dat de twee gemarkeerde punten
        $[AB]$ precies op zijn derdepunten verdelen.
  \item Pas de methode aan om een \omterm{def:g6:lines:objects}{lijnstuk} in $5$ gelijke delen te verdelen, en
        daarna om $\frac35$ van een gegeven \omterm{def:g6:lines:objects}{lijnstuk} te construeren.
  \item Construeer het punt $M$ van $[AB]$ met $\frac{AM}{MB} = \frac23$ (dat wil
        zeggen $M$ op $\frac25$ van de weg vanaf $A$). Hoeveel gelijke stappen op de
        straal, en welk punt verbind je met $B$?
  \item Waarom is deze constructie beter dan meten met een lat met maatstrepen?
        Denk aan een \omterm{def:g6:lines:objects}{lijnstuk} met lengte $\sqrt2$ (de diagonaal van een
        eenheidsvierkant, irrationaal volgens \cref{pb:g9:sqrt:1}): wat zou meten
        geven, en wat geeft Thales?
\end{enumerate}

\medskip
\noindent\textbf{Deel II --- De camera obscura.}
Prik met een speld een gaatje in één \omterm{def:g5:solids:def}{zijvlak} van een gesloten doos: op het
overstaande \omterm{def:g5:solids:def}{zijvlak} verschijnt een omgekeerd beeld van de wereld. Een punt van een
voorwerp, het gaatje en het beeldpunt liggen op één \omterm{def:g6:lines:objects}{lijn} --- licht gaat rechtdoor
--- zodat het voorwerp (hoogte $H$, op afstand $D$ vóór het gaatje) en zijn beeld
(hoogte $h$, op de achterwand op diepte $d$ achter het gaatje) in de vlinderstand
van Thales staan.
\begin{enumerate}[resume]
  \item Leid de formule van de camera obscura, $h = H \times \frac dD$, af uit
        Thales in de vlinderstand rond het gaatje.
  \item Een boom van $12$ m staat $30$ m van een schoenendoos met diepte $20$ cm.
        Hoe hoog is zijn beeld, en in welke stand?
  \item Richt de doos op de Zon: op $d = 1$ m achter het gaatje is het beeld van de
        Zon een schijf met een \omterm{def:g4:geometry:circle}{diameter} van ongeveer $9$ mm. Bereken met de afstand
        van de Zon, $D = 1.5 \times 10^{11}$ m, haar \omterm{def:g4:geometry:circle}{diameter} in \omterm{def:g9:fractions:scientific}{wetenschappelijke
        notatie} (de notatie van \cref{pb:g9:fractions:1} aan het werk). Vergelijk met
        de werkelijke waarde, $1.39 \times 10^9$ m.
  \item De Maan: \omterm{def:g4:geometry:circle}{diameter} $3.5 \times 10^6$ m, afstand $3.8 \times 10^8$ m. Bereken
        de verhouding $\frac{\text{diameter}}{\text{afstand}}$ voor de Maan en voor
        de Zon. Welk gelukkig toeval laten de twee getallen zien --- en welk
        spectaculair verschijnsel maakt het mogelijk?
  \item In één of twee zinnen: waarom staat het beeld in de camera obscura op zijn
        kop, en waarom wordt het beeld helderder maar onscherper als je het gaatje
        vergroot?
\end{enumerate}

\medskip
\noindent\textbf{Deel III --- Het trapezium van de drie gemiddelden.}
$ABCD$ is het trapezium van \cref{exo:g9:thales:9}: $(AB) \parallel (CD)$,
$AB = 4$, $CD = 6$, met de diagonalen die elkaar in $O$ snijden en
$\frac{OA}{OC} = \frac{OB}{OD} = \frac{4}{6} = \frac23$.
\begin{enumerate}[resume]
  \item Oefening op de vlinderstand: twee \omterm{def:g6:lines:objects}{lijnen} snijden elkaar in $O$ met
        $(MN) \parallel (PQ)$ in vlinderstand; $OM = 3$, $OP = 7.5$, $MN = 4$ en
        $OQ = 6$. Bereken $PQ$ en $ON$.
  \item Leg uit hoe Thales, toegepast met verhouding $\frac12$, de
        middenparallelstelling van \cref{thm:g8:midpoints:direct} als bijzonder
        geval bevat.
  \item Teken de \omterm{def:g6:lines:perp}{parallel} met de basissen door $O$; die ontmoet $[AD]$ in $E$.
        Bereken $EO$ door in de driehoek $ACD$ te werken (merk op dat
        $\frac{AO}{AC} = \frac25$).
  \item Volgens de spiegelberekening in de driehoek $BDC$ heeft het stuk $OF$ (met
        $F$ op $[BC]$) dezelfde lengte. Besluit dat
        \[
        EF = 2 \times \frac{4 \times 6}{4 + 6} = 4.8 :
        \]
        de \omterm{def:g6:lines:perp}{parallel} door het snijpunt van de diagonalen meet het \emph{\omterm{pb:g8:speed:1}{harmonisch
        gemiddelde}} van de basissen --- het gemiddelde van de heen-en-terugrit uit
        \cref{pb:g8:speed:1}, dat weer opduikt in de zuivere meetkunde.
  \item Het grote slot: bereken de drie klassieke gemiddelden van de basissen $4$
        en $6$ --- het rekenkundige $\frac{4+6}{2}$, het meetkundige
        $\sqrt{4 \times 6}$ en het harmonische $\frac{2 \times 4 \times 6}{4 + 6}$
        --- en zet ze op volgorde. Elk woont in het trapezium: het rekenkundige
        gemiddelde is de middenparallel die de middens van de benen verbindt
        (\cref{pb:g7:areas:1}), het \omterm{pb:g8:speed:1}{harmonisch gemiddelde} is jouw \omterm{def:g6:lines:objects}{lijnstuk} $EF$, en
        het \omterm{pb:g8:cosine:1}{meetkundig gemiddelde} is de \omterm{def:g6:lines:perp}{parallel} die het trapezium in twee
        \emph{gelijkvormige} trapezia verdeelt (ga zijn waarde na met de
        rekenmachine; het bewijs is een prachtige extra uitdaging). Waar is de
        ongelijkheid tussen de drie in dit boek bewezen
        (\cref{pb:g8:cosine:1}, \cref{pb:g8:speed:1})?
\end{enumerate}
\end{problem}
